\par\medskip\pdfbookmark[1]{Zu V.}{nexus-note-V}\centerline{Zu V.}\smallskip
Die zu der Formel III hinzugefügte Bemerkung giebt den Weg an, auf welchem der Verf. zur Bestimmung der Anzahl $k$ der in dem Complex $C$ enthaltenen Formen gelangt ist. Aus geometrischen Betrachtungen (vergl. I. art. 5 und II. art. 4) ergiebt sich, dass der Grenzwerth, welchem sich der Quotient
\[
\frac{(1)+(2)+(3)+\ldots+(m)}m
\]
mit unbegrenzt wachsendem $m$ nähert, d. h. die mittlere Anzahl der Darstellungen einer unbestimmten positiven ganzen Zahl
\[
=k\frac\pi{\surd D}
\]
ist; ein zweiter Ausdruck für denselben Grenzwerth lässt sich auf verschiedene Arten aus der Natur der im Vorhergehenden bestimmten Anzahl $(m)=\epsilon f(m)$ der Darstellungen der Zahl $m$ ableiten. Der zu diesem Zweck von dem Verf. zunächst eingeschlagene Weg scheint nach den vorhandenen Bruchstücken (I. artt. 7, 8; III und IV) folgender gewesen zu sein.

Ist $\theta(m)$ irgend eine Function der positiven ganzen Zahl $m$, und $p$ irgend eine Primzahl, so kann man aus $\theta(m)$ immer eine neue Function $\theta'(m)$ ableiten, deren Werth unabhängig davon ist, ob und wie oft $p$ als Factor in $m$ enthalten ist, und welche für alle durch $p$ nicht theilbaren Zahlen $m$ mit $\theta(m)$ übereinstimmt; eine solche Function erhält man, wenn man $\theta'(m)=\theta\left(\frac m{p^\pi}\right)$ setzt, wo $p^\pi$ die höchste in $m$ aufgehende Potenz von $p$ bedeutet; und man kann sagen, dass die Function $\theta'(m)$ aus $\theta(m)$ durch Elimination der Primzahl $p$ entsteht. Bildet man auf diese Weise aus $f(m)$ eine neue Function $f'(m)$ durch Elimination der Primzahl $2$, aus dieser die Function $f''(m)$ durch Elimination von $3$ u. s. f., so wird jede folgende dieser Functionen einen regelmässigern Verlauf haben, als die vorhergehenden; eliminirt man eine Primzahl nach der andern, wie sie ihrer Grösse nach auf einander folgen, so wird eine solche Function
\SourceLeaf{305}
$\theta(m)$ für unendlich viele Werthe von $m$ den Werth $f(1)=1$ haben, und namentlich für alle diejenigen Werthe von $m$, welche kleiner sind als die zuletzt eliminirte Primzahl. Durch unendliche Fortsetzung dieses Processes nähert man sich immer mehr der Function $f^\infty(m)$, welche für alle Werthe von $m$ den Werth $1$ hat, und deren mittlerer Werth folglich ebenfalls $=1$ ist. Gelingt es nun den mittlern Werth irgend einer Function $\theta(m)$ durch denjenigen der nächstfolgenden $\theta'(m)$ auszudrücken, so wird man auch den mittlern Werth der Function $f(m)$ durch eine unendliche Kette von Operationen finden können.

Ist $p$ die Primzahl, durch deren Elimination $\theta'(m)$ aus $\theta(m)$ entsteht, so ist $\theta(m)=\theta'(m)f(p^\pi)$, wenn $p^\pi$ wieder die höchste in $m$ aufgehende Potenz von $p$ bedeutet. Für den Fall, dass $p$ nicht in $2D$ aufgeht, findet man hieraus leicht, dass
\[
\theta'(m)=\theta(mp)-\left(\frac{-D}p\right)\theta(m)
\]
ist; setzt man zur Abkürzung
\[
\begin{aligned}
\Theta(m)&=\theta(1)+\theta(2)+\ldots+\theta(m)\\
\Theta'(m)&=\theta'(1)+\theta'(2)+\ldots+\theta(m)
\end{aligned}
\]
so ergiebt sich
\[
\Theta'(mp)=\Theta(mp)-\left(\frac{-D}p\right)\Theta(m)
\]
und hieraus, wenn man mit $\omega,\omega'$ resp. die mittlern Werthe der Functionen $\theta(m),\theta'(m)$ bezeichnet,
\[
\omega=\frac{\omega'}{1-\left(\frac{-D}p\right)\frac1p}
\]

Wenn aber die Primzahl $p$ in $2D$ aufgeht, so findet zwar zwischen den Functionen $\theta(m)$ und $\theta'(m)$ im Allgemeinen keine so einfache Beziehung mehr Statt; indessen ergiebt sich auf ähnliche Art leicht, dass in diesem Fall $\omega=\omega'$ ist. Ein anderer Weg, die Beziehung zwischen $\omega$ und $\omega'$ in beiden Fällen abzuleiten, ist folgender. Setzt man
\[
\vartheta(m)=\Sigma\theta(\mu)
\]
wo das Summenzeichen sich auf alle Zahlen $\mu$ bezieht, die nicht durch $p$ theilbar und ausserdem nicht grösser als $m$ sind, und bezeichnet man mit $m',m'',m'''\ldots$ resp. die grössten in $\frac mp,\frac{m'}p,\frac{m''}p\ldots$ enthaltenen ganzen Zahlen, so ist
\[
\begin{aligned}
\Theta(m)&=\vartheta(m)+\vartheta(m')f(p)+\vartheta(m'')f(pp)+\vartheta(m''')f(p^3)+\ldots\\
\Theta'(m)&=\vartheta(m)+\vartheta(m')+\vartheta(m'')+\vartheta(m''')+\ldots
\end{aligned}
\]
und hieraus folgt
\[
\frac\omega{\omega'}=(1-\tfrac1p)\{1+\frac{f(p)}p+\frac{f(pp)}{pp}+\frac{f(p^3)}{p^8}+\ldots\}
\]
was mit dem oben gefundenen Resultat übereinstimmt (vergl. die Note zu III und IV).

Der mittlere Werth der Function $f(m)$ ist daher gleich dem unendlichen Product
\[
\Pi\frac1{1-\left(\frac{-D}p\right)\frac1p}
\]
\SourceLeaf{306}
in welchem $p$ alle in $2D$ nicht aufgehenden Primzahlen durchlaufen muss, und hieraus folgt
\[
k=\frac{\epsilon\surd D}\pi\Pi\frac1{1-\left(\frac{-D}p\right)\frac1p}
\]

Hinsichtlich der Strenge dieser Deduction bleibt aber ein Bedenken übrig, welches sich auf die Methode bezieht, den mittlern Werth der Function $f(m)$ durch successive Elimination aller Primzahlen zu bestimmen; denn wenn es auch einleuchtet, dass der Werth der durch Elimination der ersten $n$ Primzahlen erhaltenen Function $f^{(n)}(m)$ mit dem der Function $f^\infty(m)=1$ übereinstimmt, so lange $m$ kleiner bleibt als die zuletzt eliminirte Primzahl, und dass also durch die Wahl eines hinreichend grossen Werthes $n$ diese Uebereinstimmung bis zu jeder vorher vorgeschriebenen Grösse der Zahl $m$ getrieben werden kann, so ist hiermit allein doch keineswegs erwiesen, dass mit unbegrenzt wachsendem $n$ der mittlere Werth der Function $f^{(n)}(m)$ sich dem mittlern Werthe der Function $f^\infty(m)$, d. h. dem Werthe $1$ unbegrenzt nähert. In welcher Weise der Verf. diese Lücke auszufüllen beabsichtigte, lässt sich aus den vorhandenen Papieren nicht mit Sicherheit erkennen; doch führt die schon oben (in der Note zu I) mitgetheilte Formel
\[
1+4\{A-\frac A3+\frac A5-\frac A7+\frac A9-\frac A{11}\ldots\}
\]
für die Anzahl der Paare von Zahlen, deren Quadratsumme den Werth $A$ nicht übertrifft, zu der Vermuthung, dass der Verf., mit Umgehung des unendlichen Productes, für den mittlern Werth der Function $f(m)$ unmittelbar die unendliche Reihe
\[
\Sigma\left(\frac{-D}n\right)\frac1n
\]
gefunden hat, in welcher $n$ der Grösse nach alle positiven ganzen Zahlen durchlaufen muss, die relative Primzahlen zu $2D$ sind. Die einfachste Art, diesen Uebergang anzudeuten, scheint die folgende zu sein.

Ist $\mu$ der grösste aller derjenigen Divisoren einer Zahl $m$, welche relative Primzahlen zu $2D$ sind, und setzt man $\theta(m)=f(\mu)$, so ist $\theta(m)$ diejenige Function, welche durch Elimination aller in $2D$ aufgehenden Primzahlen aus $f(m)$ entsteht, und deren mittlerer Werth nach dem Obigen mit demjenigen der Function $f(m)$ übereinstimmt. Da nun $\theta(m)=\Sigma\left(\frac{-D}n\right)$ ist, wo $n$ alle Divisoren von $\mu$, d. h. alle diejenigen Divisoren von $m$ durchläuft, welche relative Primzahlen zu $2D$ sind, so ergiebt sich die der obigen analoge Formel
\[
\Theta(m)=\theta(1)+\theta(2)+\ldots+\theta(m)=\Sigma\left(\frac{-D}n\right)\frac mn
\]
wo in der Summe rechter Hand der Buchstabe $n$ alle relativen Primzahlen zu $2D$ durchläuft, und von dem Quotienten $\frac mn$ immer nur die grösste in ihm enthaltene ganze Zahl beizubehalten ist. Ordnet man die Glieder dieser Reihe so, dass die Zahlen $n$ ihrer Grösse nach wachsend auf einander folgen, so nimmt der Factor $\frac mn$ fortwährend ab oder doch wenigstens nie zu, und die Reihe bricht ab, sobald $n>m$ wird. Ausserdem ergiebt sich aus dem Fundamentaltheorem in der Theorie der quadratischen Reste und aus der Verallgemeinerung desselben, dass die Summe von je $\varphi(4D)$ auf einander folgenden Werthen des Factors $\left(\frac{-D}n\right)$ verschwindet, woraus folgt, dass die Summe von noch so vielen auf einander folgenden Werthen
\SourceLeaf{307}
desselben ihrem absoluten Werth nach die endliche, nur von dem Determinant $D$ abhängige Grösse $\Delta=\varphi(2D)$ niemals übertrifft. Verbindet man diese beiden Bemerkungen mit einander, so findet man leicht, dass die Summe aller auf das Glied $\left(\frac{-D}n\right)\frac mn$ folgenden Glieder absolut genommen kleiner als $\Delta\frac mn$ ist, und dasss folglich der Quotient $\Theta(m):m$ bei unendlich wachsendem $m$ die in der angegebenen Art geordnete, convergirende unendliche Reihe
\[
\Sigma\left(\frac{-D}n\right)\frac1n
\]
zum Grenzwerthe hat. Nachdem so der gemeinschaftliche mittlere Werth der Functionen $\theta(m)$ und $f(m)$ gefunden ist, erhält man unmittelbar
\[
k=\frac{\epsilon\surd D}\pi\Sigma\left(\frac{-D}n\right)\frac1n
\]

Es verdient noch bemerkt zu werden, dass die Artikel 6 und 8 der Abhandlung II auf eine in mancher Beziehung einfachere und auch leicht auszuführende Behandlungsweise des Problems hindeuten, bei welcher nur die Darstellungen ungerader oder sogar nur solcher Zahlen betrachtet werden, die relative Primzahlen zu $2D$ sind.
\BookRule{35mm}
